The bounded dynamic results are conditional but proved: a supplied certificate gives a unitary retained interpreter run and exact undo under matching encodings (Theorem~\ref{thm:tB-sound}); Tier~A is included as a static clock schedule, without constructing a dynamic certificate (Theorem~\ref{thm:tB-A}). The exact disposal criterion requires the entire interpreter register to factor in one fixed pure state for every data input from a fixed pure initial machine state (Theorem~\ref{thm:tB-clean}). Counterexample~\ref{ex:tB-delay} restores the target at unequal completion ticks but leaves a timing record whose disposal decoheres the data. Theorem~\ref{thm:tB-cost} charges instruction-table, journal and stack scans and all supplied action, Boolean and clock circuits at every tick; its $T$-tick total is conditional on the stated circuits and access bounds. These results supply neither a general source compiler nor a certificate for an unspecified implementation.

The open-point and restriction remarks~\ref{problem:remaining}, \ref{rem:r3-restrictions} and \ref{rem:sk} delimit the following questions.
\begin{enumerate}
\item Compiling general dynamic source mark/rewind into the certified bounded interpreter remains open (Open problem~\ref{prob:tB-compiler}). The missing construction must handle mutable mark stacks, data-dependent repeated visits, general recursive calls and an inverse stopped at a coherent mark identifier, and prove the local encodings, matching routing, disjoint ranges and reachable-domain obligations after every keep. A supplied unitary schedule or static finite menu does not establish those obligations. No source-level sufficient condition for the collective factorization of Theorem~\ref{thm:tB-clean} under arbitrary intervening gates is supplied; equal padded durations do not clear the journal or stack. The obstruction
is familiar from quantum Turing machines whose computations stop after
different numbers of steps: retained records can distinguish the stopping
times and prevent the intended interference when those records are
discarded. Tier~A fixes its schedule in the syntax; Tier~B supplies an
explicit bound and padding but must still satisfy the collective disposal
test. Counterexample~\ref{ex:tB-delay} exhibits the timing-record
obstruction, and Remark~\ref{rem:tB-halting} states the checked literature
context and its limitations.
\item Theorem~\ref{thm:resA-resources} settles resource counts for the prescribed unsimplified expansion: each enclosing mark doubles body work, and this dependence is tight. These lower counts apply only to that expansion and give no lower bound for optimal equivalent circuits. Theorem~\ref{thm:resA-delayed} proves delayed erasure for the stated local evaluators with preserved basis inputs, exact basis outputs and a product private factor on each basis input, scratch-free XOR copies, and no intervening changes to retained values. It does not give a transformation for arbitrary intervening quantum gates. Keeping an inner record and deleting its immediate inverse preserves the full retained operator exactly when the commutator condition of Theorem~\ref{thm:resA-deletion} holds; on a specified input subspace that theorem supplies the weaker exact test. The five-qubit Counterexample~\ref{ex:resA-quantum} shows that inner cleanliness alone is insufficient. General optimal time--space trade-offs and unrestricted delayed-erasure transformations remain open. The conditional multiple-deletion result, Corollary~\ref{cor:resA-flat}, gives no intermediate time--space continuum.
\item Source allocation inside a mark remains excluded. A reservation translation would require a proof that the reserved factors remain zero before their first use; equality of initialized embeddings does not provide an inverse allocation on arbitrary states. The bounded deferred-observation lemma does not admit observed measurement within marked reversible computations.
\item The four-qubit calculation uses explicit Boolean copying atoms. Its equality is not asserted for the full-snapshot translation of two source read instructions, whose additional configuration records require a separate calculation.
\item General actual-phase approximation certificates remain an assumption for additional matrices beyond the effective determinant-one one-qubit rotations with a fixed number of controls in Theorem~\ref{eqv:rotations}. That theorem charges the supplied matrix-entry procedure and uses exhaustive certified search; an efficient general angle evaluator or polynomial classical synthesis time is not proved here.
\item No theorem here infers a sufficient input-independent termination horizon from the instruction list, establishes an unbounded stopping simulation, or proves correctness of an unspecified software executor. Within the settled model, position restoration and auxiliary reset remain additional conditions rather than consequences of unitary soundness.
\end{enumerate}
